\chapter{Implementation and refactoring}

\section{Engineering principles}

Treat implementation as a sequence of small, reviewable patches. Feature changes add behavior; refactoring preserves observable behavior while altering internal structure. Dependency upgrades and database migrations introduce different compatibility and rollback risks. An agent can produce a patch faster than a reviewer can validate it, so controlling change size is part of the engineering design.

\section{Implementation approach}

Establish a baseline test result and inspect the target call path. Implement one acceptance criterion, run focused tests and inspect the diff. For a migration, prefer an expand-and-contract sequence with compatibility checks, backups and rollback rehearsal. For legacy code, first add characterization tests; preserve public contracts while extracting a dependency. For cross-repository work, version APIs and coordinate integration explicitly.

\section{Verification and limitations}

Reject unrelated code churn, disabled tests, unexpected dependency changes and modifications to protected paths. A patch must pass focused and full verification before review; deployment remains a separate decision.

\section{Chapter conclusion}

Apply the stated verification criteria before expanding the scope of this workflow. Record deviations, evidence and unresolved risks in the task record.
